% Arif Mert Şengel - 2024501012
% Options for packages loaded elsewhere
\PassOptionsToPackage{unicode}{hyperref}
\PassOptionsToPackage{hyphens}{url}
\PassOptionsToPackage{dvipsnames,svgnames,x11names}{xcolor}
\documentclass[
  10pt,
  a4paper,
  titlepage]{article}
\usepackage{xcolor}
\usepackage[margin=22mm]{geometry}
\usepackage{amsmath,amssymb}
\setcounter{secnumdepth}{5}
\usepackage{iftex}
\ifPDFTeX
  \usepackage[T1]{fontenc}
  \usepackage[utf8]{inputenc}
  \usepackage{textcomp} % provide euro and other symbols
\else % if luatex or xetex
  \usepackage{unicode-math} % this also loads fontspec
  \defaultfontfeatures{Scale=MatchLowercase}
  \defaultfontfeatures[\rmfamily]{Ligatures=TeX,Scale=1}
\fi
\usepackage{lmodern}
\ifPDFTeX\else
  % xetex/luatex font selection
\fi
% Use upquote if available, for straight quotes in verbatim environments
\IfFileExists{upquote.sty}{\usepackage{upquote}}{}
\IfFileExists{microtype.sty}{% use microtype if available
  \usepackage[]{microtype}
  \UseMicrotypeSet[protrusion]{basicmath} % disable protrusion for tt fonts
}{}
\makeatletter
\@ifundefined{KOMAClassName}{% if non-KOMA class
  \IfFileExists{parskip.sty}{%
    \usepackage{parskip}
  }{% else
    \setlength{\parindent}{0pt}
    \setlength{\parskip}{6pt plus 2pt minus 1pt}}
}{% if KOMA class
  \KOMAoptions{parskip=half}}
\makeatother
\usepackage{listings}
\newcommand{\passthrough}[1]{#1}
\lstset{defaultdialect=[5.3]Lua}
\lstset{defaultdialect=[x86masm]Assembler}
\usepackage{longtable,booktabs,array}
\usepackage{caption}
\newcounter{none} % for unnumbered tables
\usepackage{calc} % for calculating minipage widths
% Correct order of tables after \paragraph or \subparagraph
\usepackage{etoolbox}
\makeatletter
\patchcmd\longtable{\par}{\if@noskipsec\mbox{}\fi\par}{}{}
\makeatother
\setlength{\emergencystretch}{3em} % prevent overfull lines
\providecommand{\tightlist}{%
  \setlength{\itemsep}{0pt}\setlength{\parskip}{0pt}}
% Arif Mert Şengel - 2024501012
\usepackage{fontspec}
% Use the TeX-bundled Latin Modern defaults rather than OS-installed fonts.
\usepackage{xcolor,listings,booktabs,longtable,array,fancyhdr,lastpage,titlesec,xurl,needspace}
\definecolor{guideink}{HTML}{18334D}
\definecolor{guideaccent}{HTML}{245D85}
\definecolor{guidepaper}{HTML}{F5F7FA}
\definecolor{guideborder}{HTML}{CBD4DF}
\definecolor{guidecomment}{HTML}{536B55}
\definecolor{guidestring}{HTML}{8C4D31}
\lstdefinelanguage{JavaScript}{
  keywords={const,let,var,function,return,for,of,in,if,else,try,catch,finally,throw,new,class,extends,import,from,export,default,yield,true,false,null,undefined},
  sensitive=true,
  morecomment=[l]{//},
  morecomment=[s]{/*}{*/},
  morestring=[b]',
  morestring=[b]",
  morestring=[b]`
}
\lstdefinelanguage{TypeScript}[]{JavaScript}{morekeywords={type,infer,never,unknown,readonly,keyof,typeof,number,string,boolean,as}}
\lstdefinelanguage{JSX}[]{JavaScript}{}
\lstdefinelanguage{TSX}[]{TypeScript}{}
\lstdefinelanguage{JavaScriptReact}[]{JavaScript}{}
\lstdefinelanguage{TypeScriptReact}[]{TypeScript}{}
\lstdefinelanguage{js}[]{JavaScript}{}
\lstdefinelanguage{jsx}[]{JavaScript}{}
\lstdefinelanguage{ts}[]{TypeScript}{}
\lstdefinelanguage{tsx}[]{TypeScript}{}
\lstdefinelanguage{text}{}
\lstdefinelanguage{Text}{}
\lstdefinelanguage{bash}{morekeywords={npm,node,npx,git,cd,cp,rm,bash,python3,tectonic},sensitive=true,morecomment=[l]{\#},morestring=[b]",morestring=[b]'}
\lstdefinelanguage{sh}[]{bash}{}
\lstset{
  basicstyle=\ttfamily\fontsize{8.2pt}{10.3pt}\selectfont,
  keywordstyle=\color{guideaccent},
  commentstyle=\color{guidecomment},
  stringstyle=\color{guidestring},
  backgroundcolor=\color{guidepaper},
  frame=single,rulecolor=\color{guideborder},framerule=0.3pt,framesep=5pt,
  breaklines=true,breakatwhitespace=false,columns=fullflexible,
  keepspaces=true,showstringspaces=false,upquote=true,
  xleftmargin=0.7em,xrightmargin=0.7em,
  aboveskip=0.7em,belowskip=0.7em,
  literate={Ç}{{\c{C}}}1 {ç}{{\c{c}}}1 {Ş}{{\c{S}}}1 {ş}{{\c{s}}}1
           {Ğ}{{\u{G}}}1 {ğ}{{\u{g}}}1 {İ}{{\.I}}1 {ı}{{\i}}1
           {Ö}{{\"O}}1 {ö}{{\"o}}1 {Ü}{{\"U}}1 {ü}{{\"u}}1
}
\setlength{\parindent}{0pt}
\setlength{\parskip}{4pt plus 1pt}
\setlength{\emergencystretch}{3em}
\setlength{\headheight}{14pt}
\setcounter{secnumdepth}{2}
\setcounter{tocdepth}{1}
\pagestyle{fancy}
\fancyhf{}
\fancyhead[L]{\small\sffamily\color{guideink}Hardrock field guide}
\fancyhead[R]{\small\sffamily MIS 233}
\fancyfoot[L]{\small Arif Mert Şengel}
\fancyfoot[R]{\small \thepage\ / \pageref*{LastPage}}
\renewcommand{\headrulewidth}{0.3pt}
\renewcommand{\footrulewidth}{0pt}
\titleformat{\section}{\normalfont\Large\bfseries\color{guideink}}{\thesection}{0.6em}{}
\titleformat{\subsection}{\normalfont\large\bfseries\color{guideink}}{\thesubsection}{0.6em}{}
\titlespacing*{\section}{0pt}{1.2em}{0.6em}
\titlespacing*{\subsection}{0pt}{1em}{0.4em}
\pretocmd{\section}{\Needspace{20\baselineskip}}{}{}
\pretocmd{\subsection}{\Needspace{15\baselineskip}}{}{}
\AtBeginEnvironment{longtable}{\small\setlength{\tabcolsep}{4pt}}
\let\guidetoc\tableofcontents
\renewcommand{\tableofcontents}{\guidetoc\clearpage}
\let\guidemaketitle\maketitle
\renewcommand{\maketitle}{\hypersetup{pageanchor=false}\guidemaketitle\hypersetup{pageanchor=true}}
\urlstyle{same}
\Urlmuskip=0mu plus 1mu
\usepackage{bookmark}
\IfFileExists{xurl.sty}{\usepackage{xurl}}{} % add URL line breaks if available
\urlstyle{same}
% fallback for those not using the hyperref driver hyperxmp:
\makeatletter
\@ifundefined{xmpquote}{\newcommand{\xmpquote}[1]{#1}}{}
\makeatother
\hypersetup{
  pdftitle={Destructuring Field Guide - Hardrock},
  pdfauthor={Arif Mert Şengel - 2024501012},
  colorlinks=true,
  linkcolor={guideink},
  filecolor={Maroon},
  citecolor={guideink},
  urlcolor={guideaccent},
  pdfcreator={LaTeX via pandoc}}

\title{Destructuring Field Guide - Hardrock}
\usepackage{etoolbox}
\makeatletter
\providecommand{\subtitle}[1]{% add subtitle to \maketitle
  \apptocmd{\@title}{\par {\large #1 \par}}{}{}
}
\makeatother
\subtitle{JavaScript, TypeScript and React - MIS 233}
\author{Arif Mert Şengel - 2024501012}
\date{21 November 2026}

\begin{document}
\maketitle

{
\setcounter{tocdepth}{1}
\tableofcontents
}
\section{About this guide}\label{about-this-guide}

This guide explores JavaScript and TypeScript destructuring through my
Hardrock playlist. The examples cover positional values, named
properties, defaults, rest, parameter patterns and React props and
children. The personal log records my reported predictions, experiments
and actual terminal results.

Bonuses claimed: Bonus 1: LaTeX; Bonus 2: Depth; Bonus 3: Real code;
Bonus 4: Git repository.

\section{My data and the example
settings}\label{my-data-and-the-example-settings}

{\def\LTcaptype{none} % do not increment counter
\begin{longtable}[]{@{}lllll@{}}
\toprule\noalign{}
Title & Artist & Duration I supplied & Seconds & My rating \\
\midrule\noalign{}
\endhead
\bottomrule\noalign{}
\endlastfoot
Siyah & Çilekeş & 4:45 & 285 & 10/10 \\
Uyan & Pickpocket & 3:51 & 231 & 7/10 \\
Yine & Kurban & 4:21 & 261 & 9/10 \\
\end{longtable}
}

Playlist: Hardrock. Listening context: ders çalışırken (while studying).
These are the values I supplied, not verified music-catalog metadata.
Saved volume/repeat settings and mutation examples below are
simulations, not additional claims about my listening history.

\section{Run}\label{run}

Use Node.js 20.19 or newer. In the folder containing package.json:

\Needspace{61.8pt}
\begin{lstlisting}[language=sh]
npm install
npm run all
npm run check
npm run all > output.txt
\end{lstlisting}

Verification environment: Node.js v26.13.0, TypeScript 5.8.3,
React/react-dom 19.3.1 and tsx 4.23.15. npm run all and npm run check
completed successfully. The recorded output.txt contains real output
from npm run all. Regenerate it after changes to the final source. A
successful tsx execution does not replace tsc type checking.

The starter strict, noUncheckedIndexedAccess and react-jsx settings are
preserved. There is no explicit any or @ts-ignore in the source. The
sole deliberate suppression is the T5 @ts-expect-error. The original
all/js/ts/check scripts are preserved; npm run depth runs the additional
JavaScript and TypeScript investigations.

\section{What destructuring does}\label{what-destructuring-does}

An array pattern takes successive values from an iterable; an object
pattern selects properties by name. A declaration introduces bindings,
while assignment destructuring updates existing targets. Neither
operation deletes the original properties. A const binding cannot be
reassigned, but objects it refers to can still be mutated.

Rest collects remaining values:
\passthrough{\lstinline![first, ...remaining]!} or
\passthrough{\lstinline!\{ rating, ...publicTrack \}!}. Spread expands a
value in an expression: \passthrough{\lstinline![...queue]!},
\passthrough{\lstinline!\{ ...song \}!} or
\passthrough{\lstinline!<section \{...rest\}>!}. They share the same
dots but perform different jobs. Formal array and object binding are
defined in the
\href{https://tc39.es/ecma262/multipage/ecmascript-language-statements-and-declarations.html\#sec-destructuring-binding-patterns}{ECMAScript
specification}.

The plain equivalents below explain these concrete examples, not every
JavaScript edge case. Array patterns also support non-array iterables.
Repeating a property access in a handwritten default expression can
trigger a getter more than once. Form 8 uses object spread for its
shallow copy, while avoiding destructuring.

\section{Field guide}\label{field-guide}

Each form is implemented in both languages. Outputs in this section are
the JS run, which matched the corresponding TS run. They are displayed
once rather than repeating an entire second run; repeated lines within a
single experiment are preserved.

\subsection{Form 1: Return the playlist duration and average
rating}\label{form-1-return-the-playlist-duration-and-average-rating}

\textbf{Syntax:}
\passthrough{\lstinline!const [playlistSeconds, averageRating] = summarizePlaylist(songs)!}

\Needspace{9\baselineskip}

\textbf{Example} (\passthrough{\lstinline!js/arrays.mjs!}):

\Needspace{82.4pt}
\begin{lstlisting}[language=JavaScript]
const songs = [
  {
    title: 'Siyah',
    artist: 'Çilekeş',
    durationSeconds: 285,
    rating: 10,
  },
  {
    title: 'Uyan',
    artist: 'Pickpocket',
    durationSeconds: 231,
    rating: 7,
  },
  {
    title: 'Yine',
    artist: 'Kurban',
    durationSeconds: 261,
    rating: 9,
  },
]

function summarizePlaylist(tracks) {
  let totalSeconds = 0
  let totalRating = 0

  for (const track of tracks) {
    totalSeconds += track.durationSeconds
    totalRating += track.rating
  }

  return [totalSeconds, totalRating / tracks.length]
}

const [playlistSeconds, averageRating] = summarizePlaylist(songs)

console.log('F1:', playlistSeconds, averageRating)
\end{lstlisting}

\Needspace{9\baselineskip}

\textbf{Output} (actual run):

\Needspace{30.9pt}
\begin{lstlisting}
F1: 777 8.666666666666666
\end{lstlisting}

\textbf{Use it when:} A summary calculation returns two related results
that I want to use separately: duration and average score.

\textbf{Watch out:} Positions determine the meaning. Renaming the two
local variables does not change the order returned by the function. This
calculation assumes a nonempty playlist; an empty playlist needs an
explicit policy for its average.

\textbf{TypeScript} (\passthrough{\lstinline!ts/arrays.ts!}): Use a
Song{[}{]} parameter and a {[}number, number{]} return type. The tuple
promises two numeric positions, rather than an ordinary array of unknown
length.

\Needspace{9\baselineskip}

\textbf{Without destructuring:}

\Needspace{51.5pt}
\begin{lstlisting}[language=JavaScript]
const result = summarizePlaylist(songs)
const playlistSeconds = result[0]
const averageRating = result[1]
\end{lstlisting}

\subsection{Form 2: Skip the current song and collect the remaining
queue}\label{form-2-skip-the-current-song-and-collect-the-remaining-queue}

\textbf{Syntax:}
\passthrough{\lstinline!const [, nextSong, ...remainingSongs] = queue!}

\Needspace{9\baselineskip}

\textbf{Example} (\passthrough{\lstinline!js/arrays.mjs!}):

\Needspace{61.8pt}
\begin{lstlisting}[language=JavaScript]
const queue = ['Siyah', 'Uyan', 'Yine']
const [, nextSong, ...remainingSongs] = queue
console.log('F2 next:', nextSong)
console.log('F2 remaining:', JSON.stringify(remainingSongs))
\end{lstlisting}

\Needspace{9\baselineskip}

\textbf{Output} (actual run):

\Needspace{41.2pt}
\begin{lstlisting}
F2 next: Uyan
F2 remaining: ["Yine"]
\end{lstlisting}

\textbf{Use it when:} The first song is already playing. I need the next
song and everything after it without modifying the original queue.

\textbf{Watch out:} The leading empty position consumes Siyah.
remainingSongs is a new array containing Yine. If the queue is too
short, nextSong can be undefined. Rest must be last and cannot have a
trailing comma.

\textbf{TypeScript} (\passthrough{\lstinline!ts/arrays.ts!}): The queue
is string{[}{]}. With noUncheckedIndexedAccess, nextSong is string
\textbar{} undefined, while remainingSongs is string{[}{]}.

\Needspace{9\baselineskip}

\textbf{Without destructuring:}

\Needspace{41.2pt}
\begin{lstlisting}[language=JavaScript]
const nextSong = queue[1]
const remainingSongs = queue.slice(2)
\end{lstlisting}

\subsection{Form 3: Swap the opening and closing
songs}\label{form-3-swap-the-opening-and-closing-songs}

\textbf{Syntax:}
\passthrough{\lstinline![openingSong, closingSong] = [closingSong, openingSong]!}

\Needspace{9\baselineskip}

\textbf{Example} (\passthrough{\lstinline!js/arrays.mjs!}):

\Needspace{72.1pt}
\begin{lstlisting}[language=JavaScript]
let openingSong = 'Siyah'
let closingSong = 'Yine'
console.log('F3 before:', openingSong, closingSong)
;[openingSong, closingSong] = [closingSong, openingSong]
console.log('F3 after:', openingSong, closingSong)
\end{lstlisting}

\Needspace{9\baselineskip}

\textbf{Output} (actual run):

\Needspace{41.2pt}
\begin{lstlisting}
F3 before: Siyah Yine
F3 after: Yine Siyah
\end{lstlisting}

\textbf{Use it when:} I want to reverse two selected positions without
losing the original opening song.

\textbf{Watch out:} Both variables must be mutable and already declared.
The right-hand array is evaluated before the assignments. The leading
semicolon prevents this array-starting statement from joining the
preceding expression.

\textbf{TypeScript} (\passthrough{\lstinline!ts/arrays.ts!}): Both
variables have type string. Swapping a string and a number would be
rejected. No new destructuring-specific type syntax is needed.

\Needspace{9\baselineskip}

\textbf{Without destructuring:}

\Needspace{51.5pt}
\begin{lstlisting}[language=JavaScript]
const previousOpening = openingSong
openingSong = closingSong
closingSong = previousOpening
\end{lstlisting}

\subsection{Form 4: Defaults for missing saved playback
settings}\label{form-4-defaults-for-missing-saved-playback-settings}

\textbf{Syntax:}
\passthrough{\lstinline!const [volume = 35, repeat = false] = savedSettings!}

\Needspace{9\baselineskip}

\textbf{Example} (\passthrough{\lstinline!js/arrays.mjs!}):

\Needspace{61.8pt}
\begin{lstlisting}[language=JavaScript]
// These are simulated saved settings, not claimed listening preferences.
const savedSettings = []
const [volume = 35, repeat = false] = savedSettings
console.log('F4:', volume, repeat)
\end{lstlisting}

\Needspace{9\baselineskip}

\textbf{Output} (actual run):

\Needspace{30.9pt}
\begin{lstlisting}
F4: 35 false
\end{lstlisting}

\textbf{Use it when:} A simulated saved-settings array may omit the
volume or repeat flag.

\textbf{Watch out:} An item default replaces undefined, including a
missing item. It does not replace null, zero or false. The settings
values are invented for this experiment, not claimed as my actual player
configuration.

\textbf{TypeScript} (\passthrough{\lstinline!ts/arrays.ts!}): The
optional tuple {[}number?, boolean?{]} describes both positions. After
these defaults, volume is number and repeat is boolean.

\Needspace{9\baselineskip}

\textbf{Without destructuring:}

\Needspace{41.2pt}
\begin{lstlisting}[language=JavaScript]
const volume = savedSettings[0] === undefined ? 35 : savedSettings[0]
const repeat = savedSettings[1] === undefined ? false : savedSettings[1]
\end{lstlisting}

\subsection{Form 5: Select playlist properties by
name}\label{form-5-select-playlist-properties-by-name}

\textbf{Syntax:}
\passthrough{\lstinline!let \{ name, activity \} = playlist!}

\Needspace{9\baselineskip}

\textbf{Example} (\passthrough{\lstinline!js/objects.mjs!}):

\Needspace{51.5pt}
\begin{lstlisting}[language=JavaScript]
const playlist = { name: 'Hardrock', activity: 'ders çalışırken' }
let { name, activity } = playlist
console.log('F5:', name, activity)
\end{lstlisting}

\Needspace{9\baselineskip}

\textbf{Output} (actual run):

\Needspace{30.9pt}
\begin{lstlisting}
F5: Hardrock ders çalışırken
\end{lstlisting}

\textbf{Use it when:} The playlist name and listening context are used
together in a local calculation or display.

\textbf{Watch out:} Property names select the values regardless of the
property order in these plain data objects. Reassigning a local string
does not update playlist. A null or undefined source still throws.

\textbf{TypeScript} (\passthrough{\lstinline!ts/objects.ts!}): The
object type declares name and activity as strings. A misspelled property
is rejected instead of silently producing an undefined local value.

\Needspace{9\baselineskip}

\textbf{Without destructuring:}

\Needspace{41.2pt}
\begin{lstlisting}[language=JavaScript]
let name = playlist.name
let activity = playlist.activity
\end{lstlisting}

\subsection{Form 6: Rename track fields for the current
context}\label{form-6-rename-track-fields-for-the-current-context}

\textbf{Syntax:}
\passthrough{\lstinline!const \{ title: trackTitle, artist: artistName \} = song!}

\Needspace{9\baselineskip}

\textbf{Example} (\passthrough{\lstinline!js/objects.mjs!}):

\Needspace{51.5pt}
\begin{lstlisting}[language=JavaScript]
const song = { title: 'Siyah', artist: 'Çilekeş' }
const { title: trackTitle, artist: artistName } = song
console.log('F6:', trackTitle, artistName)
\end{lstlisting}

\Needspace{9\baselineskip}

\textbf{Output} (actual run):

\Needspace{30.9pt}
\begin{lstlisting}
F6: Siyah Çilekeş
\end{lstlisting}

\textbf{Use it when:} I want descriptive local names, especially when
playlist and song information are displayed together.

\textbf{Watch out:} The original property is on the left of each colon;
the local variable is on the right. The pattern creates trackTitle and
artistName, not local title and artist variables.

\textbf{TypeScript} (\passthrough{\lstinline!ts/objects.ts!}): The colon
inside the pattern still means rename. Any explicit object-type
annotation goes after the entire pattern.

\Needspace{9\baselineskip}

\textbf{Without destructuring:}

\Needspace{41.2pt}
\begin{lstlisting}[language=JavaScript]
const trackTitle = song.title
const artistName = song.artist
\end{lstlisting}

\subsection{Form 7: Default a playlist
note}\label{form-7-default-a-playlist-note}

\textbf{Syntax:}
\passthrough{\lstinline!const \{ note = 'No playlist note' \} = playlist!}

\Needspace{9\baselineskip}

\textbf{Example} (\passthrough{\lstinline!js/objects.mjs!}):

\Needspace{51.5pt}
\begin{lstlisting}[language=JavaScript]
const playlist = { name: 'Hardrock' }
const { note = 'No playlist note' } = playlist
console.log('F7:', note)
\end{lstlisting}

\Needspace{9\baselineskip}

\textbf{Output} (actual run):

\Needspace{30.9pt}
\begin{lstlisting}
F7: No playlist note
\end{lstlisting}

\textbf{Use it when:} A playlist can have an optional note, and the
display needs text when none was supplied.

\textbf{Watch out:} The default applies to undefined, including an
absent note. It neither writes a note back into playlist nor replaces an
explicit null.

\textbf{TypeScript} (\passthrough{\lstinline!ts/objects.ts!}): note?:
string permits the absent property. Its local default removes undefined
from the result type. If null were also allowed, it would remain
possible.

\Needspace{9\baselineskip}

\textbf{Without destructuring:}

\Needspace{30.9pt}
\begin{lstlisting}[language=JavaScript]
const note = playlist.note === undefined ? 'No playlist note' : playlist.note
\end{lstlisting}

\subsection{Form 8: Separate my rating from public track
information}\label{form-8-separate-my-rating-from-public-track-information}

\textbf{Syntax:}
\passthrough{\lstinline!const \{ rating, ...publicTrack \} = song!}

\Needspace{9\baselineskip}

\textbf{Example} (\passthrough{\lstinline!js/objects.mjs!}):

\Needspace{72.1pt}
\begin{lstlisting}[language=JavaScript]
const song = { title: 'Siyah', artist: 'Çilekeş', durationSeconds: 285, rating: 10 }
// Select my personal rating separately from the public track information.
const { rating, ...publicTrack } = song
console.log('F8 rating:', rating)
console.log('F8 public:', JSON.stringify(publicTrack))
\end{lstlisting}

\Needspace{9\baselineskip}

\textbf{Output} (actual run):

\Needspace{41.2pt}
\begin{lstlisting}
F8 rating: 10
F8 public: {"title":"Siyah","artist":"Çilekeş","durationSeconds":285}
\end{lstlisting}

\textbf{Use it when:} I want to handle my personal score separately
while keeping title, artist and duration together.

\textbf{Watch out:} publicTrack is a new outer object that excludes
rating. Object rest copies remaining own enumerable properties and does
not deep-clone nested objects.

\textbf{TypeScript} (\passthrough{\lstinline!ts/objects.ts!}): Song
describes the original fields. The inferred rest type excludes rating,
so publicTrack.rating is a compile error.

\Needspace{9\baselineskip}

\textbf{Without destructuring:}

\Needspace{51.5pt}
\begin{lstlisting}[language=JavaScript]
const rating = song.rating
const publicTrack = { ...song }
delete publicTrack.rating
\end{lstlisting}

\subsection{Form 9: Select a nested listening context and the first
track}\label{form-9-select-a-nested-listening-context-and-the-first-track}

\textbf{Syntax:}
\passthrough{\lstinline!const \{ listening: \{ activity \}, tracks: [firstTrack] \} = playlist!}

\Needspace{9\baselineskip}

\textbf{Example} (\passthrough{\lstinline!js/objects.mjs!}):

\Needspace{82.4pt}
\begin{lstlisting}[language=JavaScript]
const playlist = {
  listening: { activity: 'ders çalışırken' },
  tracks: [{ title: 'Siyah' }, { title: 'Uyan' }, { title: 'Yine' }],
}
const { listening: { activity }, tracks: [firstTrack] } = playlist
console.log('F9:', activity, firstTrack?.title)
\end{lstlisting}

\Needspace{9\baselineskip}

\textbf{Output} (actual run):

\Needspace{30.9pt}
\begin{lstlisting}
F9: ders çalışırken Siyah
\end{lstlisting}

\textbf{Use it when:} A playlist has both nested metadata and an array
of tracks, and I only need two small pieces of that structure.

\textbf{Watch out:} This binds activity and firstTrack, not listening or
tracks. A missing listening object fails before any optional access on
firstTrack. The ?. in the log handles an empty tracks array; the pattern
itself is not optional chaining.

\textbf{TypeScript} (\passthrough{\lstinline!ts/objects.ts!}): tracks is
an array of objects. firstTrack can be undefined, so both versions use
firstTrack?.title in the log. The listening parent is required by the
type.

\Needspace{9\baselineskip}

\textbf{Without destructuring:}

\Needspace{41.2pt}
\begin{lstlisting}[language=JavaScript]
const activity = playlist.listening.activity
const firstTrack = playlist.tracks[0]
\end{lstlisting}

\subsection{Form 10: Choose a song field using a
variable}\label{form-10-choose-a-song-field-using-a-variable}

\textbf{Syntax:}
\passthrough{\lstinline!const \{ [key]: value \} = song!}

\Needspace{9\baselineskip}

\textbf{Example} (\passthrough{\lstinline!js/objects.mjs!}):

\Needspace{82.4pt}
\begin{lstlisting}[language=JavaScript]
const song = { title: 'Siyah', durationSeconds: 285, rating: 10 }
function readSongField(key) {
  const { [key]: value } = song
  return value
}
console.log('F10 rating:', readSongField('rating'))
console.log('F10 duration:', readSongField('durationSeconds'))
\end{lstlisting}

\Needspace{9\baselineskip}

\textbf{Output} (actual run):

\Needspace{41.2pt}
\begin{lstlisting}
F10 rating: 10
F10 duration: 285
\end{lstlisting}

\textbf{Use it when:} The same helper selects either a rating or a
duration according to the requested field name.

\textbf{Watch out:} The brackets evaluate key. They do not request a
property literally named key. Selecting different fields can also
produce different value types.

\textbf{TypeScript} (\passthrough{\lstinline!ts/objects.ts!}): keyof
typeof song restricts the key to real fields. The helper returns string
\textbar{} number because title is also a permitted field.

\Needspace{9\baselineskip}

\textbf{Without destructuring:}

\Needspace{30.9pt}
\begin{lstlisting}[language=JavaScript]
const value = song[key]
\end{lstlisting}

\subsection{Form 11: Destructure playback options in
parameters}\label{form-11-destructure-playback-options-in-parameters}

\textbf{Syntax:}
\passthrough{\lstinline!function playbackOptions(\{ playlistName = 'Hardrock', volume = 35, repeat = false \} = \{\})!}

\Needspace{9\baselineskip}

\textbf{Example} (\passthrough{\lstinline!js/params.mjs!}):

\Needspace{82.4pt}
\begin{lstlisting}[language=JavaScript]
// Volume and repeat are simulated settings for this example.
function playbackOptions({ playlistName = 'Hardrock', volume = 35, repeat = false } = {}) {
  return { playlistName, volume, repeat }
}
console.log('F11 omitted:', JSON.stringify(playbackOptions()))
console.log('F11 supplied:', JSON.stringify(playbackOptions({ volume: 60 })))
function toSeconds([minutes, seconds]) {
  return minutes * 60 + seconds
}
console.log('F11 array parameter:', toSeconds([4, 45]))
\end{lstlisting}

\Needspace{9\baselineskip}

\textbf{Output} (actual run):

\Needspace{51.5pt}
\begin{lstlisting}
F11 omitted: {"playlistName":"Hardrock","volume":35,"repeat":false}
F11 supplied: {"playlistName":"Hardrock","volume":60,"repeat":false}
F11 array parameter: 285
\end{lstlisting}

\textbf{Use it when:} An options object makes named playback settings
readable and permits callers to leave out settings.

\textbf{Watch out:} The outer = \{\} handles an omitted or undefined
argument. The inner defaults handle missing or undefined fields. Neither
layer replaces an explicit null argument. The settings are simulated
examples.

\textbf{TypeScript} (\passthrough{\lstinline!ts/params.ts!}):
PlaybackInput goes after the whole pattern. Its fields are optional so
\{\} is a valid default. The additional array-parameter example uses a
{[}number, number{]} tuple.

\Needspace{9\baselineskip}

\textbf{Without destructuring:}

\Needspace{82.4pt}
\begin{lstlisting}[language=JavaScript]
function playbackOptions(options = {}) {
  const playlistName = options.playlistName === undefined ? 'Hardrock' : options.playlistName
  const volume = options.volume === undefined ? 35 : options.volume
  const repeat = options.repeat === undefined ? false : options.repeat
  return { playlistName, volume, repeat }
}
function toSeconds(pair) {
  const minutes = pair[0]
  const seconds = pair[1]
  return minutes * 60 + seconds
}
\end{lstlisting}

\subsection{Form 12: Destructure entries in rating and artist
loops}\label{form-12-destructure-entries-in-rating-and-artist-loops}

\textbf{Syntax:}
\passthrough{\lstinline!for (const [title, rating] of Object.entries(ratings))!}

\Needspace{9\baselineskip}

\textbf{Example} (\passthrough{\lstinline!js/params.mjs!}):

\Needspace{82.4pt}
\begin{lstlisting}[language=JavaScript]
const ratings = { Siyah: 10, Uyan: 7, Yine: 9 }
for (const [title, rating] of Object.entries(ratings)) {
  console.log('F12 object:', title, rating)
}
const artists = new Map([['Siyah', 'Çilekeş'], ['Uyan', 'Pickpocket'], ['Yine', 'Kurban']])
for (const [title, artist] of artists) {
  console.log('F12 Map:', title, artist)
}
\end{lstlisting}

\Needspace{9\baselineskip}

\textbf{Output} (actual run):

\Needspace{82.4pt}
\begin{lstlisting}
F12 object: Siyah 10
F12 object: Uyan 7
F12 object: Yine 9
F12 Map: Siyah Çilekeş
F12 Map: Uyan Pickpocket
F12 Map: Yine Kurban
\end{lstlisting}

\textbf{Use it when:} Ratings and artist names are naturally associated
with track titles, and a loop needs both parts of each entry.

\textbf{Watch out:} Object.entries returns own enumerable string-keyed
entries. A Map already yields key/value pairs. for...in yields
property-name strings, not these entry pairs, and a plain object is not
itself iterable.

\textbf{TypeScript} (\passthrough{\lstinline!ts/params.ts!}):
Record\textless string, number\textgreater{} describes the rating
object, and Map\textless string, string\textgreater{} describes the
artist entries. Object.entries does not preserve a precise union of the
literal title keys here.

\Needspace{9\baselineskip}

\textbf{Without destructuring:}

\Needspace{82.4pt}
\begin{lstlisting}[language=JavaScript]
for (const entry of Object.entries(ratings)) {
  const title = entry[0]
  const rating = entry[1]
  console.log('F12 object:', title, rating)
}
for (const entry of artists) {
  const title = entry[0]
  const artist = entry[1]
  console.log('F12 Map:', title, artist)
}
\end{lstlisting}

\section{Pitfalls}\label{pitfalls}

The five experiments are in js/pitfalls.mjs. Errors are intentionally
caught so the complete project can continue running.

\subsection{Pitfall 1: An absent source
throws}\label{pitfall-1-an-absent-source-throws}

\Needspace{82.4pt}
\begin{lstlisting}[language=JavaScript]
for (const input of [undefined, null]) {
  try {
    const { title } = input
    console.log(title)
  } catch (error) {
    console.log(`Pitfall 1 ${String(input)}:`, `${error.name}: ${error.message}`)
  }
}
\end{lstlisting}

\Needspace{9\baselineskip}

\textbf{Real output:}

\Needspace{41.2pt}
\begin{lstlisting}
Pitfall 1 undefined: TypeError: Cannot destructure property 'title' of 'input' as it is undefined.
Pitfall 1 null: TypeError: Cannot destructure property 'title' of 'input' as it is null.
\end{lstlisting}

\textbf{Why:} An object pattern cannot read properties from null or
undefined. Validate the input, or use input ?? \{\} if both values
should mean an empty object. A parameter = \{\} default handles
undefined, but still does not replace null.

\subsection{Pitfall 2: Defaults replace undefined, not null or
zero}\label{pitfall-2-defaults-replace-undefined-not-null-or-zero}

\Needspace{82.4pt}
\begin{lstlisting}[language=JavaScript]
const cases = [
  ['missing', {}],
  ['undefined', { volume: undefined }],
  ['null', { volume: null }],
  ['zero', { volume: 0 }],
]
for (const [label, settings] of cases) {
  const { volume = 35 } = settings
  console.log(`Pitfall 2 ${label}:`, volume)
}
\end{lstlisting}

\Needspace{9\baselineskip}

\textbf{Real output:}

\Needspace{61.8pt}
\begin{lstlisting}
Pitfall 2 missing: 35
Pitfall 2 undefined: 35
Pitfall 2 null: null
Pitfall 2 zero: 0
\end{lstlisting}

\textbf{Why:} Both the absent volume property and the explicit undefined
use 35. Explicit null and 0 remain unchanged. The four labelled lines
are all preserved; a repeated numeric result is meaningful evidence.

\subsection{Pitfall 3: Object assignment at statement start needs
parentheses}\label{pitfall-3-object-assignment-at-statement-start-needs-parentheses}

\Needspace{82.4pt}
\begin{lstlisting}[language=JavaScript]
try {
  new Function('let title; { title } = { title: "Siyah" };')()
} catch (error) {
  console.log('Pitfall 3 invalid:', `${error.name}: ${error.message}`)
}
let title
;({ title } = { title: 'Siyah' })
console.log('Pitfall 3 fixed:', title)
\end{lstlisting}

\Needspace{9\baselineskip}

\textbf{Real output:}

\Needspace{41.2pt}
\begin{lstlisting}
Pitfall 3 invalid: SyntaxError: Unexpected token '='
Pitfall 3 fixed: Siyah
\end{lstlisting}

\textbf{Why:} The parser treats a statement beginning with \{ as a
block. Parentheses force the expression interpretation. The invalid
attempt is parsed with new Function so its SyntaxError can be caught;
placing invalid syntax directly in this file would stop the entire file
from parsing.

\subsection{Pitfall 4: Rest is a shallow
copy}\label{pitfall-4-rest-is-a-shallow-copy}

\Needspace{82.4pt}
\begin{lstlisting}[language=JavaScript]
// This mutation is an experiment, not a change to my recorded song rating.
const song = { title: 'Siyah', stats: { rating: 10 } }
const { title, ...copy } = song
copy.stats.rating = 7
console.log('Pitfall 4 original rating:', song.stats.rating)
console.log('Pitfall 4 different outer object:', copy !== song)
console.log('Pitfall 4 shared stats:', copy.stats === song.stats)
\end{lstlisting}

\Needspace{9\baselineskip}

\textbf{Real output:}

\Needspace{51.5pt}
\begin{lstlisting}
Pitfall 4 original rating: 7
Pitfall 4 different outer object: true
Pitfall 4 shared stats: true
\end{lstlisting}

\textbf{Why:} The copy is a different outer object, but its stats
property still points to the original nested object. Changing
copy.stats.rating therefore changes song.stats.rating. This isolated
experiment does not alter the source dataset used by Form 1.

\subsection{Pitfall 5: Object names and array positions select
differently}\label{pitfall-5-object-names-and-array-positions-select-differently}

\Needspace{72.1pt}
\begin{lstlisting}[language=JavaScript]
const song = { title: 'Siyah', artist: 'Çilekeş' }
const { artist, title } = song
const [arrayArtist, arrayTitle] = [song.title, song.artist]
console.log('Pitfall 5 object:', artist, title)
console.log('Pitfall 5 array:', arrayArtist, arrayTitle)
\end{lstlisting}

\Needspace{9\baselineskip}

\textbf{Real output:}

\Needspace{41.2pt}
\begin{lstlisting}
Pitfall 5 object: Çilekeş Siyah
Pitfall 5 array: Siyah Çilekeş
\end{lstlisting}

\textbf{Why:} The object pattern selects artist by its property name
even though artist is written first. The array pattern selects the first
value, Siyah, for arrayArtist despite that misleading variable name.
Array destructuring does not infer meaning from local names.

\section{TypeScript topics T1-T5}\label{typescript-topics-t1t5}

All five are implemented in ts/types.ts. Types check the contracts; they
do not change the runtime meaning of a binding pattern.

\subsection{T1: A colon in a pattern
renames}\label{t1-a-colon-in-a-pattern-renames}

\Needspace{61.8pt}
\begin{lstlisting}[language=TypeScript]
const source = { rating: 10 }
const { rating: string } = source
const { rating }: { rating: number } = source
console.log('T1:', string, typeof string, rating)
\end{lstlisting}

\Needspace{9\baselineskip}

\textbf{Real output:}

\Needspace{30.9pt}
\begin{lstlisting}
T1: 10 number 10
\end{lstlisting}

The variable named string receives the number 10. It is not a type
annotation. The correct annotation is after the whole pattern:
\passthrough{\lstinline!const \{ rating \}: \{ rating: number \} = source!}.
\href{https://www.typescriptlang.org/docs/handbook/variable-declarations.html\#destructuring}{TypeScript
destructuring documentation}.

\subsection{T2: Tuple return types and as
const}\label{t2-tuple-return-types-and-as-const}

\Needspace{82.4pt}
\begin{lstlisting}[language=TypeScript]
function favouriteSong(): [number, string] {
  return [10, 'Siyah']
}
const [rating, title] = favouriteSong()
const fixed = [10, 'Siyah'] as const
const [fixedRating, fixedTitle] = fixed
const ten: 10 = fixedRating
const siyah: 'Siyah' = fixedTitle
console.log('T2:', rating, title, ten, siyah)
\end{lstlisting}

\Needspace{9\baselineskip}

\textbf{Real output:}

\Needspace{30.9pt}
\begin{lstlisting}
T2: 10 Siyah 10 Siyah
\end{lstlisting}

favouriteSong returns {[}number, string{]}, so its first item is numeric
and its second is a string. The literal as const value has type readonly
{[}10, "Siyah"{]}, and its bindings retain those literal types. as const
does not freeze or deep-clone a runtime object. The separate compiler
probe rejects fixed.push(9).
\href{https://www.typescriptlang.org/docs/handbook/release-notes/typescript-3-4.html\#const-assertions}{TypeScript
const assertions}.

\subsection{T3: An ordinary array may be
empty}\label{t3-an-ordinary-array-may-be-empty}

\Needspace{72.1pt}
\begin{lstlisting}[language=TypeScript]
const queue: string[] = []
const [first] = queue // string | undefined
const [safeFirst = 'Siyah'] = queue // string
const message = first === undefined ? 'No song selected' : first.toUpperCase()
console.log('T3:', first, safeFirst, message)
\end{lstlisting}

\Needspace{9\baselineskip}

\textbf{Real output:}

\Needspace{30.9pt}
\begin{lstlisting}
T3: undefined Siyah No song selected
\end{lstlisting}

Under noUncheckedIndexedAccess, first is string \textbar{} undefined. A
default produces a definite string, while an explicit undefined check
narrows first before toUpperCase. A non-null assertion would merely hide
the missing-value possibility.
\href{https://www.typescriptlang.org/tsconfig/noUncheckedIndexedAccess.html}{Compiler
option reference}.

\subsection{T4: Optional properties and defaulted
bindings}\label{t4-optional-properties-and-defaulted-bindings}

\Needspace{82.4pt}
\begin{lstlisting}[language=TypeScript]
type Settings = { volume?: number; note?: string | null }
const settings: Settings = { note: null }
const { volume: rawVolume } = settings // number | undefined
const { volume = 35, note = 'No playlist note' } = settings
const readyVolume: number = volume
const nullableNote: string | null = note
console.log('T4:', rawVolume, readyVolume, nullableNote)
\end{lstlisting}

\Needspace{9\baselineskip}

\textbf{Real output:}

\Needspace{30.9pt}
\begin{lstlisting}
T4: undefined 35 null
\end{lstlisting}

rawVolume is number \textbar{} undefined. volume is number after its
default. note remains string \textbar{} null because null does not
activate a default. Optional and nullable are different conditions.

\subsection{T5: An intentional compile error and its
fix}\label{t5-an-intentional-compile-error-and-its-fix}

\Needspace{82.4pt}
\begin{lstlisting}[language=TypeScript]
const song = { title: 'Siyah' }
// @ts-expect-error T5: Property 'titel' does not exist; the property is title.
const { titel } = song
const { title } = song
console.log('T5 wrong:', titel)
console.log('T5 fixed:', title)
\end{lstlisting}

\Needspace{9\baselineskip}

\textbf{Real output:}

\Needspace{41.2pt}
\begin{lstlisting}
T5 wrong: undefined
T5 fixed: Siyah
\end{lstlisting}

The misspelling produces
\passthrough{\lstinline!error TS2339: Property 'titel' does not exist on type '\{ title: string; \}'.!}
The separate probe observed this message, recorded in
compiler-errors.txt. The sole @ts-expect-error marks that one intended
error in the real source. At runtime the missing property becomes
undefined; selecting title fixes the code. An expected-error directive
is not general permission to suppress other mistakes.

\section{Props and children}\label{props-and-children}

A function component receives one props object. children is one field of
that object; it is not a second independent package. JSX passed to
another named prop, such as header, remains that named prop.
React\textquotesingle s
\href{https://react.dev/learn/passing-props-to-a-component}{props guide}
describes passing values and nested content.

The JS examples below import renderToStaticMarkup from react-dom/server.
js/props.jsx also imports format from node:util to print the captured
diagnostic. Matching typed implementations are in ts/props.tsx, using:

\Needspace{72.1pt}
\begin{lstlisting}[language=TSX]
import type { ReactNode, PropsWithChildren, ComponentProps } from 'react'
import { renderToStaticMarkup } from 'react-dom/server'
type PlaylistProps = { name: string; children?: ReactNode }
type FrameProps = { children?: ReactNode }
type ShellProps = { name: string } & ComponentProps<'section'>
\end{lstlisting}

The output is non-interactive server-rendered HTML, not a browser
playback test. Event handlers do not execute in this render.
\href{https://react.dev/reference/react-dom/server/renderToStaticMarkup}{renderToStaticMarkup
reference}.

\subsection{P1: Props is one object}\label{p1-props-is-one-object}

\Needspace{82.4pt}
\begin{lstlisting}[language=JSX]
function PlaylistA(props) {
  return <section><h2>{props.name}</h2>{props.children}</section>
}
function PlaylistB({ name, children }) {
  return <section><h2>{name}</h2>{children}</section>
}
console.log('P1 A:', renderToStaticMarkup(
  <PlaylistA name="Hardrock"><p>ders çalışırken</p></PlaylistA>))
console.log('P1 B:', renderToStaticMarkup(
  <PlaylistB name="Hardrock"><p>ders çalışırken</p></PlaylistB>))
\end{lstlisting}

\Needspace{9\baselineskip}

\textbf{Real rendered output:}

\Needspace{41.2pt}
\begin{lstlisting}
P1 A: <section><h2>Hardrock</h2><p>ders çalışırken</p></section>
P1 B: <section><h2>Hardrock</h2><p>ders çalışırken</p></section>
\end{lstlisting}

PlaylistA reads name and children from props in its body. PlaylistB
selects those fields in its parameter. The identical markup demonstrates
the same contract. The latter combines Form 5 with Form 11: an object
binding pattern in a function parameter. The TS signatures place
PlaylistProps after the whole parameter or pattern.

\subsection{P2: children is a prop}\label{p2-children-is-a-prop}

\Needspace{82.4pt}
\begin{lstlisting}[language=JSX]
console.log('P2 nested:', renderToStaticMarkup(
  <PlaylistB name="Hardrock">Siyah</PlaylistB>))
console.log('P2 attribute:', renderToStaticMarkup(
  <PlaylistB name="Hardrock" children="Siyah" />))
const supplied = { children: 'Uyan' }
console.log('P2 collision:', renderToStaticMarkup(
  <PlaylistB name="Hardrock" {...supplied}>Siyah</PlaylistB>))
\end{lstlisting}

\Needspace{9\baselineskip}

\textbf{Real rendered output:}

\Needspace{51.5pt}
\begin{lstlisting}
P2 nested: <section><h2>Hardrock</h2>Siyah</section>
P2 attribute: <section><h2>Hardrock</h2>Siyah</section>
P2 collision: <section><h2>Hardrock</h2>Siyah</section>
\end{lstlisting}

The nested Siyah text and the explicit children attribute produce the
same HTML. The collision example spreads children="Uyan" and also
supplies nested Siyah text; nested content replaces the spread value.
This illustrates precedence. A caller should normally supply one
unambiguous children source.

\subsection{P3: Renderable children and a function
contract}\label{p3-renderable-children-and-a-function-contract}

\Needspace{82.4pt}
\begin{lstlisting}[language=JSX]
function Frame({ children }) {
  return <div>{children}</div>
}
const cases = [
  ['string', 'Siyah'],
  ['number', 10],
  ['element', <em>Çilekeş</em>],
  ['several', [<b key="siyah">Siyah</b>, <i key="uyan">Uyan</i>]],
  ['undefined', undefined],
  ['null', null],
  ['false', false],
]
for (const [label, child] of cases) {
  console.log(`P3 ${label}:`, renderToStaticMarkup(<Frame>{child}</Frame>))
}
const renderTrackCount = count => <strong>{count} tracks</strong>
// RAW FUNCTION START
// This deliberately invalid JS attempt records React's actual diagnostic.
const diagnostics = []
const originalError = console.error
let rawMarkup
try {
  console.error = (...args) => diagnostics.push(format(...args))
  rawMarkup = renderToStaticMarkup(<Frame>{renderTrackCount}</Frame>)
} finally {
  console.error = originalError
}
console.log('P3 raw function:', rawMarkup)
for (const diagnostic of diagnostics) console.log('P3 diagnostic:', diagnostic)
// RAW FUNCTION END
// A function child works when the receiving component explicitly calls it.
function PlaylistSummary({ count, children }) {
  return <div>{children(count)}</div>
}
console.log('P3 called function:', renderToStaticMarkup(
  <PlaylistSummary count={3}>{renderTrackCount}</PlaylistSummary>))
\end{lstlisting}

\Needspace{9\baselineskip}

\textbf{Real rendered output:}

\Needspace{82.4pt}
\begin{lstlisting}
P3 string: <div>Siyah</div>
P3 number: <div>10</div>
P3 element: <div><em>Çilekeş</em></div>
P3 several: <div><b>Siyah</b><i>Uyan</i></div>
P3 undefined: <div></div>
P3 null: <div></div>
P3 false: <div></div>
P3 raw function: <div></div>
P3 diagnostic: Functions are not valid as a React child. This may happen if you return renderTrackCount instead of <renderTrackCount /> from render. Or maybe you meant to call this function rather than return it.
P3 called function: <div><strong>3 tracks</strong></div>
\end{lstlisting}

The string, number 10, one element and array of keyed elements appear in
the HTML. undefined, null and false produce no content inside the div. A
raw function is not called by React: the JS experiment records the
actual diagnostic and an empty div. Its temporary console.error
replacement is restored in finally and the diagnostic is printed into
stdout. PlaylistSummary instead explicitly accepts a function and calls
children(count), producing the 3 tracks element. Its TS contract is
\passthrough{\lstinline!\{ count: number; children: (count: number) => ReactNode \}!}.
An ordinary Frame with children?: ReactNode rejects an uncalled
function. The negative compiler probe demonstrates that rejection
without adding a second suppression to the main project.

\subsection{P4: Collect children separately and forward
rest}\label{p4-collect-children-separately-and-forward-rest}

\Needspace{82.4pt}
\begin{lstlisting}[language=JSX]
function FixedStyle({ name, children, ...rest }) {
  return <section {...rest} className="music"><h2>{name}</h2>{children}</section>
}
function CallerStyle({ name, children, ...rest }) {
  return <section className="music" {...rest}><h2>{name}</h2>{children}</section>
}
console.log('P4 explicit last:', renderToStaticMarkup(
  <FixedStyle name="Hardrock" id="study-list" className="study">Siyah</FixedStyle>))
console.log('P4 spread last:', renderToStaticMarkup(
  <CallerStyle name="Hardrock" id="study-list" className="study">Siyah</CallerStyle>))
\end{lstlisting}

\Needspace{9\baselineskip}

\textbf{Real rendered output:}

\Needspace{41.2pt}
\begin{lstlisting}
P4 explicit last: <section id="study-list" class="music"><h2>Hardrock</h2>Siyah</section>
P4 spread last: <section class="study" id="study-list"><h2>Hardrock</h2>Siyah</section>
\end{lstlisting}

The parameter uses object rest to collect unhandled fields, then JSX
spread forwards them to section. FixedStyle writes className last, so
music wins. CallerStyle spreads last, so the supplied study value wins.
id is forwarded in both. name is consumed as the heading, not forwarded.
The type combines the custom name with native section props; the code
should forward only intended DOM-compatible values. The attribute
ordering in the strings can differ without changing which class value
won.

\subsection{P5: Named slots or children-only
composition}\label{p5-named-slots-or-children-only-composition}

\Needspace{82.4pt}
\begin{lstlisting}[language=JSX]
function PlayerWithSlots({ header, footer, children }) {
  return <section><header>{header}</header><main>{children}</main><footer>{footer}</footer></section>
}
function PlayerWithChildren({ children }) {
  return <section>{children}</section>
}
const trackList = <ol><li>Siyah</li><li>Uyan</li><li>Yine</li></ol>
console.log('P5 named slots:', renderToStaticMarkup(
  <PlayerWithSlots header={<h3>Hardrock</h3>} footer={<small>3 tracks</small>}>
    {trackList}
  </PlayerWithSlots>))
console.log('P5 children only:', renderToStaticMarkup(
  <PlayerWithChildren>
    <header><h3>Hardrock</h3></header><main>{trackList}</main>
    <footer><small>3 tracks</small></footer>
  </PlayerWithChildren>))
\end{lstlisting}

\Needspace{9\baselineskip}

\textbf{Real rendered output:}

\Needspace{41.2pt}
\begin{lstlisting}
P5 named slots: <section><header><h3>Hardrock</h3></header><main><ol><li>Siyah</li><li>Uyan</li><li>Yine</li></ol></main><footer><small>3 tracks</small></footer></section>
P5 children only: <section><header><h3>Hardrock</h3></header><main><ol><li>Siyah</li><li>Uyan</li><li>Yine</li></ol></main><footer><small>3 tracks</small></footer></section>
\end{lstlisting}

Both approaches produce the same section/header/main/footer markup.
PlayerWithSlots owns that structure and lets the caller supply header,
footer and body contents. This suits a consistent player layout.
PlayerWithChildren lets the caller choose the entire inner structure and
ordering, which suits open-ended composition. In the slot version header
and footer are named props and the ol is children. In the children-only
version the header, main and footer are all children.

\subsection{P6: Undefined triggers a prop default; null does
not}\label{p6-undefined-triggers-a-prop-default-null-does-not}

\Needspace{82.4pt}
\begin{lstlisting}[language=JSX]
function PlaylistNote({ note = 'No listening note' }) {
  return <p>{note}</p>
}
console.log('P6 missing:', renderToStaticMarkup(<PlaylistNote />))
console.log('P6 undefined:', renderToStaticMarkup(<PlaylistNote note={undefined} />))
console.log('P6 null:', renderToStaticMarkup(<PlaylistNote note={null} />))
console.log('P6 empty string:', renderToStaticMarkup(<PlaylistNote note="" />))
\end{lstlisting}

\Needspace{9\baselineskip}

\textbf{Real rendered output:}

\Needspace{61.8pt}
\begin{lstlisting}
P6 missing: <p>No listening note</p>
P6 undefined: <p>No listening note</p>
P6 null: <p></p>
P6 empty string: <p></p>
\end{lstlisting}

The missing and explicitly undefined note use the fallback. null remains
null and produces an empty p. The empty string also stays unchanged and
renders no text. This is Form 7 inside a parameter pattern. The TS type
deliberately allows note?: string \textbar{} null so the null experiment
is legal.

\subsection{P7: Three ways to describe
props}\label{p7-three-ways-to-describe-props}

\Needspace{82.4pt}
\begin{lstlisting}[language=JSX]
function RatingBox({ children }) {
  return <aside>{children}</aside>
}
function ListeningBox({ label, children }) {
  return <aside><strong>{label}</strong>{children}</aside>
}
function QueueButton({ children, ...rest }) {
  return <button {...rest}>{children}</button>
}
console.log('P7 ReactNode:', renderToStaticMarkup(<RatingBox>{10}</RatingBox>))
console.log('P7 PropsWithChildren:', renderToStaticMarkup(
  <ListeningBox label="Studying">Hardrock</ListeningBox>))
console.log('P7 ComponentProps:', renderToStaticMarkup(
  <QueueButton type="button" disabled aria-label="Queue Siyah">Queue</QueueButton>))
\end{lstlisting}

\Needspace{9\baselineskip}

\textbf{Real rendered output:}

\Needspace{51.5pt}
\begin{lstlisting}
P7 ReactNode: <aside>10</aside>
P7 PropsWithChildren: <aside><strong>Studying</strong>Hardrock</aside>
P7 ComponentProps: <button type="button" disabled="" aria-label="Queue Siyah">Queue</button>
\end{lstlisting}

children?: ReactNode explicitly allows omitted children and renderable
values such as numbers, strings and elements. It rejects an ordinary
uncalled function and an arbitrary plain object.
PropsWithChildren\textless\{ label: string \}\textgreater{} adds
optional children while preserving the required label.
ComponentProps\textless\textquotesingle button\textquotesingle\textgreater{}
obtains the native button props, including type, disabled, onClick and
children; it does not forward or render them by itself. The spread in
QueueButton performs the forwarding.
\href{https://react.dev/learn/typescript}{React TypeScript guide}.

Typed implementations:

\Needspace{82.4pt}
\begin{lstlisting}[language=TSX]
function RatingBox({ children }: { children?: ReactNode }) {
  return <aside>{children}</aside>
}
function ListeningBox({ label, children }: PropsWithChildren<{ label: string }>) {
  return <aside><strong>{label}</strong>{children}</aside>
}
function QueueButton({ children, ...rest }: ComponentProps<'button'>) {
  return <button {...rest}>{children}</button>
}
\end{lstlisting}

The separate compiler experiment rejected an uncalled function for
RatingBox, a missing required label for ListeningBox, href for
QueueButton, and a string disabled value. These attempts are saved as
text in compiler-probe.txt, not left as errors in the passing source.
The actual messages include:

\Needspace{72.1pt}
\begin{lstlisting}
error TS2322: Type '() => string' is not assignable to type 'ReactNode'.
error TS2741: Property 'label' is missing in type '{ children: string; }' but required in type '{ label: string; }'.
error TS2322: Type '{ children: string; href: string; }' is not assignable to type 'IntrinsicAttributes & ClassAttributes<HTMLButtonElement> & ButtonHTMLAttributes<HTMLButtonElement>'.
  Property 'href' does not exist on type 'IntrinsicAttributes & ClassAttributes<HTMLButtonElement> & ButtonHTMLAttributes<HTMLButtonElement>'. Did you mean 'ref'?
error TS2322: Type 'string' is not assignable to type 'boolean | undefined'.
\end{lstlisting}

Props are read-only inputs to a component. Mutating props is not a way
to update a parent. The key used on the list elements is React
bookkeeping, not an ordinary readable props.key; pass another identifier
prop if a component needs it.
\href{https://react.dev/warnings/special-props}{React key warning}.

\section{JavaScript versus
TypeScript}\label{javascript-versus-typescript}

The first two columns describe the same runtime patterns. The errors
below are actual results from the deliberately invalid compiler probe;
they are not unresolved errors in npm run check. Inference often
supplies types without an explicit annotation. Union bars are displayed
as \textbar{} in this table to keep its columns intact; the untouched
ASCII \textbar{} messages are in compiler-errors.txt.

{\def\LTcaptype{none} % do not increment counter
\begin{longtable}[]{@{}
  >{\raggedright\arraybackslash}p{(\linewidth - 6\tabcolsep) * \real{0.0600}}
  >{\raggedright\arraybackslash}p{(\linewidth - 6\tabcolsep) * \real{0.2400}}
  >{\raggedright\arraybackslash}p{(\linewidth - 6\tabcolsep) * \real{0.3000}}
  >{\raggedright\arraybackslash}p{(\linewidth - 6\tabcolsep) * \real{0.4000}}@{}}
\toprule\noalign{}
\begin{minipage}[b]{\linewidth}\raggedright
Form
\end{minipage} & \begin{minipage}[b]{\linewidth}\raggedright
What stays the same
\end{minipage} & \begin{minipage}[b]{\linewidth}\raggedright
Type information used or needed
\end{minipage} & \begin{minipage}[b]{\linewidth}\raggedright
Actual error in the wrong version
\end{minipage} \\
\midrule\noalign{}
\endhead
\bottomrule\noalign{}
\endlastfoot
1 & Position selects each returned result & A tuple return, such as
{[}number, number{]} & TS2322:
\passthrough{\lstinline!Type 'string | number | undefined' is not assignable to type 'number'.!} \\
2 & Skip one item and collect the rest & string{[}{]} and a possible
undefined indexed item & TS18048:
\passthrough{\lstinline!'first' is possibly 'undefined'.!} \\
3 & RHS values are evaluated before swapping & Compatible types for both
existing variables & TS2322:
\passthrough{\lstinline!Type 'number' is not assignable to type 'string'.!} \\
4 & Only undefined activates the default & Optional tuple positions and
compatible default values & TS2322:
\passthrough{\lstinline!Type 'number | "unknown"' is not assignable to type 'number'.!} \\
5 & Select properties by name & The declared/inferred object fields &
TS2339:
\passthrough{\lstinline!Property 'titel' does not exist on type '\{ title: string; \}'.!} \\
6 & A colon renames a local binding & Annotation outside the pattern, if
needed & TS2304: \passthrough{\lstinline!Cannot find name 'artist'.!} \\
7 & A field default does not replace null & Optional fields; a default
removes undefined & TS2322:
\passthrough{\lstinline!Type 'number | undefined' is not assignable to type 'number'.!} \\
8 & Selected fields are excluded from rest & Inferred rest type without
rating & TS2339:
\passthrough{\lstinline!Property 'rating' does not exist on type '\{ title: string; \}'.!} \\
9 & Nested patterns follow the source shape & Required/guarded parents,
possible missing array item & TS2339:
\passthrough{\lstinline!Property 'activity' does not exist on type '\{ activity: string; \} | undefined'.!} \\
10 & {[}key{]} evaluates a property name & keyof typeof song for a
known-key helper & TS2537:
\passthrough{\lstinline!Type '\{ title: string; rating: number; \}' has no matching index signature for type 'string'.!} \\
11 & Parameter and property defaults are separate & An input type after
the entire pattern & TS2741:
\passthrough{\lstinline!Property 'playlistName' is missing in type '\{\}' but required in type '\{ playlistName: string; \}'.!} \\
12 & Object.entries and Map yield entry pairs & Record and Map key/value
types & TS2488:
\passthrough{\lstinline!Type '\{ Siyah: number; \}' must have a '[Symbol.iterator]()' method that returns an iterator.!} \\
\end{longtable}
}

keyof is a type-level union of permitted property names, not a runtime
transform.
\href{https://www.typescriptlang.org/docs/handbook/2/keyof-types.html}{TypeScript
keyof reference}.

\section{Reproducing the negative compiler
experiment}\label{reproducing-the-negative-compiler-experiment}

compiler-probe.txt contains the deliberately invalid examples for the
table and P7. compiler-errors.txt contains the actual diagnostics from
that separate compiler run. To reproduce separately:

\Needspace{61.8pt}
\begin{lstlisting}[language=sh]
cp compiler-probe.txt _compiler_probe.tsx
node node_modules/typescript/bin/tsc --noEmit --strict --noUncheckedIndexedAccess --skipLibCheck --target ES2022 --module NodeNext --moduleResolution NodeNext --jsx react-jsx _compiler_probe.tsx --pretty false > compiler-errors.txt
rm _compiler_probe.tsx
npm run check
\end{lstlisting}

The explicit probe command should fail. The generated
\_compiler\_probe.tsx must then be removed; the normal project check
must pass. Source-file line numbers and wording can vary with compiler
versions, so record what your own run actually prints.

\section{Bonus 1: LaTeX}\label{bonus-1-latex}

The complete field guide is supplied as latex/main.tex and
latex/main.pdf. It includes all twelve form entries, the five pitfalls,
T1-T5, P1-P7, the comparison table, the depth investigations and the
personal log. The title/name fields appear on the title page. The
contents page links to sections, source/output blocks use listings, and
the tables are actual longtable tables rather than screenshots.

The source was generated from this README with Pandoc 3.14 and compiled
with Tectonic 0.17.0. The conversion helpers are in scripts/. To
regenerate the TeX source with Pandoc installed, then compile the PDF
with Tectonic installed:

\Needspace{41.2pt}
\begin{lstlisting}[language=sh]
python3 scripts/make_latex.py
bash scripts/build_pdf.sh
\end{lstlisting}

The compiler command executed by the build script is:

\Needspace{30.9pt}
\begin{lstlisting}[language=sh]
tectonic -X compile --keep-logs --outdir latex latex/main.tex
\end{lstlisting}

The build also accepts -\/-pandoc /path/to/pandoc and the TECTONIC\_BIN
environment variable for a compiler outside PATH. The included main.tex
is self-contained, so Pandoc is needed only when regenerating the source
after editing README. Tectonic runs the required LaTeX passes and
produces the PDF. No shell-escape or external syntax-highlighting
process is required. The PDF was checked for clipped text, missing
characters and overfull boxes; build details are recorded in
latex/build-report.txt.

\section{Bonus 2: Depth}\label{bonus-2-depth}

The investigations below use the Hardrock data to examine iterable
access, evaluation order, copying boundaries, unusual binding locations
and type-level extraction. Each runtime result comes from a real run of
the included files. No execution-time benchmark or general performance
claim is inferred from these small experiments.

Sources: js/depth.mjs and ts/depth.tsx. Run both with:

\Needspace{30.9pt}
\begin{lstlisting}[language=sh]
npm run depth
\end{lstlisting}

The full captured output is depth-output.txt. The TypeScript part is
included in the ordinary strict npm run check. The deliberate negative
cases are kept separately as text.

\subsection{D1: Strings, Set and Map are iterable
values}\label{d1-strings-set-and-map-are-iterable-values}

\Needspace{82.4pt}
\begin{lstlisting}[language=JavaScript]
const [firstLetter, ...remainingLetters] = 'Çilekeş'
console.log('D1 string:', firstLetter, remainingLetters.join(''))
const [firstUnique, ...remainingUnique] = new Set(['Siyah', 'Uyan', 'Siyah', 'Yine'])
console.log('D1 Set:', firstUnique, JSON.stringify(remainingUnique))
const [[title, rating]] = new Map([['Siyah', 10]])
console.log('D1 Map:', title, rating)
\end{lstlisting}

\Needspace{9\baselineskip}

\textbf{Real output:}

\Needspace{51.5pt}
\begin{lstlisting}
D1 string: Ç ilekeş
D1 Set: Siyah ["Uyan","Yine"]
D1 Map: Siyah 10
\end{lstlisting}

\textbf{What the experiment shows:} The string yields Ç followed by the
remaining characters. Set discards the second Siyah while keeping
insertion order. Map yields an entry array, so the nested pattern gets
its key and value. These examples establish that array patterns accept
iterable values, not only Array instances. The following trace
experiment tests the protocol directly.

\subsection{D2: Observe iterator, next and return
calls}\label{d2-observe-iterator-next-and-return-calls}

\Needspace{82.4pt}
\begin{lstlisting}[language=JavaScript]
function tracedQueue(label) {
  const titles = ['Siyah', 'Uyan', 'Yine']
  return {
    [Symbol.iterator]() {
      let index = 0
      console.log(`${label} iterator`)
      return {
        next() {
          const done = index >= titles.length
          console.log(`${label} next:`, index, done ? 'done' : 'value')
          return { value: titles[index++], done }
        },
        return() {
          console.log(`${label} return`)
          return { done: true }
        },
      }
    },
  }
}
const [first] = tracedQueue('D2 one')
console.log('D2 one result:', first)
const [, second] = tracedQueue('D2 skip')
console.log('D2 skip result:', second)
const [head, ...tail] = tracedQueue('D2 rest')
console.log('D2 rest result:', head, JSON.stringify(tail))
const [] = tracedQueue('D2 empty')
const { 0: numericProperty } = tracedQueue('D2 object')
console.log('D2 object property:', numericProperty)
\end{lstlisting}

\Needspace{9\baselineskip}

\textbf{Real output:}

\Needspace{82.4pt}
\begin{lstlisting}
D2 one iterator
D2 one next: 0 value
D2 one return
D2 one result: Siyah
D2 skip iterator
D2 skip next: 0 value
D2 skip next: 1 value
D2 skip return
D2 skip result: Uyan
D2 rest iterator
D2 rest next: 0 value
D2 rest next: 1 value
D2 rest next: 2 value
D2 rest next: 3 done
D2 rest result: Siyah ["Uyan","Yine"]
D2 empty iterator
D2 empty return
D2 object property: undefined
\end{lstlisting}

\textbf{What the experiment shows:} Taking one song calls the iterator
method once, next once, then return to close the unfinished iterator.
Skipping a position still calls next, so the second-song example
consumes two values before closing. Rest requests all three values and
one final done result; it does not call return after ordinary
exhaustion. The empty {[}{]} pattern still acquires and closes an
iterator without calling next. By contrast, the object pattern asks for
property 0 and never invokes the iterator method; this object has no
such property, so the value is undefined. These are different access
mechanisms, not interchangeable indexing syntax. The
\href{https://tc39.es/ecma262/multipage/abstract-operations.html\#sec-getiterator}{GetIterator
and IteratorClose algorithms} explain the traced operations.

\subsection{D3: Compare early closure with complete generator
consumption}\label{d3-compare-early-closure-with-complete-generator-consumption}

\Needspace{82.4pt}
\begin{lstlisting}[language=JavaScript]
function* playQueue(label) {
  try {
    for (const title of ['Siyah', 'Uyan', 'Yine']) {
      console.log(`${label} requested:`, title)
      yield title
    }
  } finally {
    console.log(`${label} closed`)
  }
}
const [first] = playQueue('D3 short')
console.log('D3 short result:', first)
const [head, ...rest] = playQueue('D3 full')
console.log('D3 full result:', head, JSON.stringify(rest))
\end{lstlisting}

\Needspace{9\baselineskip}

\textbf{Real output:}

\Needspace{82.4pt}
\begin{lstlisting}
D3 short requested: Siyah
D3 short closed
D3 short result: Siyah
D3 full requested: Siyah
D3 full requested: Uyan
D3 full requested: Yine
D3 full closed
D3 full result: Siyah ["Uyan","Yine"]
\end{lstlisting}

\textbf{What the experiment shows:} The short pattern requests only
Siyah, and the generator finally block runs before the result is
printed. Adding rest requests Siyah, Uyan and Yine before finally runs.
In the short case, closing the unfinished generator triggers finally; in
the full case, finally also runs when the generator reaches normal
completion. A finally log alone therefore does not prove that return was
called. This is why the explicit protocol trace in D2 is useful. Rest
also means full consumption: collecting an infinite queue would never
finish.

\subsection{D4: Lazy defaults, dependencies and getter
order}\label{d4-lazy-defaults-dependencies-and-getter-order}

\Needspace{82.4pt}
\begin{lstlisting}[language=JavaScript]
let calls = 0
function defaultVolume() {
  calls += 1
  return 35
}
const [volume = defaultVolume(), quiet = volume - 10,
  note = defaultVolume(), muted = defaultVolume()] = [undefined, undefined, null, 0]
console.log('D4 defaults:', volume, quiet, note, muted, calls)
try {
  const [first = later, later = 'Yine'] = []
  console.log(first, later)
} catch (error) {
  console.log('D4 later binding:', `${error.name}: ${error.message}`)
}
const events = []
const settings = {
  get volume() { events.push('read volume'); return undefined },
  get label() { events.push('read label'); return 'Hardrock' },
}
function chooseKey() { events.push('choose key'); return 'volume' }
function fallback() { events.push('default'); return 35 }
const { [chooseKey()]: selectedVolume = fallback(), ...rest } = settings
console.log('D4 evaluation order:', JSON.stringify(events))
console.log('D4 selected:', selectedVolume, JSON.stringify(rest))
\end{lstlisting}

\Needspace{9\baselineskip}

\textbf{Real output:}

\Needspace{61.8pt}
\begin{lstlisting}
D4 defaults: 35 25 null 0 1
D4 later binding: ReferenceError: Cannot access 'later' before initialization
D4 evaluation order: ["choose key","read volume","default","read label"]
D4 selected: 35 {"label":"Hardrock"}
\end{lstlisting}

\textbf{What the experiment shows:} The four bindings print 35, 25, null
and 0 with exactly one defaultVolume call. An earlier binding is
available to the next default, while null and zero prevent their
initializers from running. Reading a later binding before initialization
throws the recorded ReferenceError. The second experiment records choose
key, read volume, default and read label in that order. The selected
getter is read before its default is chosen; rest then reads the other
enumerable field. A default expression can have observable side effects,
so moving it elsewhere or repeatedly reading a getter is not always an
equivalent refactoring. The
\href{https://tc39.es/ecma262/multipage/syntax-directed-operations.html\#sec-runtime-semantics-iteratorbindinginitialization}{binding
initialization algorithms} distinguish extraction from conditional
default evaluation.

\subsection{D5: Object rest, inheritance, non-enumerable fields and
symbols}\label{d5-object-rest-inheritance-non-enumerable-fields-and-symbols}

\Needspace{82.4pt}
\begin{lstlisting}[language=JavaScript]
const sourceTag = Symbol('source')
const song = Object.create({ playlistName: 'Hardrock' })
song.title = 'Siyah'
song.rating = 10
song[sourceTag] = 'student list'
Object.defineProperty(song, 'internalNote', { value: 'local only', enumerable: false })
const { playlistName, ...publicTrack } = song
console.log('D5 inherited lookup:', playlistName)
console.log('D5 string keys:', JSON.stringify(Object.keys(publicTrack)))
console.log('D5 symbol keys:', JSON.stringify(Object.getOwnPropertySymbols(publicTrack).map(key => key.description)))
console.log('D5 hidden copied:', Object.hasOwn(publicTrack, 'internalNote'))
console.log('D5 JSON:', JSON.stringify(publicTrack))
\end{lstlisting}

\Needspace{9\baselineskip}

\textbf{Real output:}

\Needspace{72.1pt}
\begin{lstlisting}
D5 inherited lookup: Hardrock
D5 string keys: ["title","rating"]
D5 symbol keys: ["source"]
D5 hidden copied: false
D5 JSON: {"title":"Siyah","rating":10}
\end{lstlisting}

\textbf{What the experiment shows:} A named property lookup can find
playlistName on the prototype. Rest does not copy that inherited field
or the non-enumerable internalNote. It does copy the own enumerable
symbol property. Object.keys shows only title and rating, but
getOwnPropertySymbols reveals source. JSON output also omits the symbol
key, so JSON.stringify alone is insufficient evidence of everything rest
copied. The experiment explicitly checks each kind of property rather
than assuming that the JSON string describes the whole object.

\subsection{D6: A loop can destructure the wrong kind of
value}\label{d6-a-loop-can-destructure-the-wrong-kind-of-value}

\Needspace{82.4pt}
\begin{lstlisting}[language=JavaScript]
const ratings = { Siyah: 10, Uyan: 7, Yine: 9 }
for (const [first, second] in ratings) {
  console.log('D6 for-in:', first, second)
}
for (const [title, rating] of Object.entries(ratings)) {
  console.log('D6 for-of:', title, rating)
}
const current = { title: '', rating: 0 }
;[current.title, current.rating] = ['Siyah', 10]
console.log('D6 assignment targets:', JSON.stringify(current))
try {
  throw new Error('No song selected')
} catch ({ message }) {
  console.log('D6 catch:', message)
}
\end{lstlisting}

\Needspace{9\baselineskip}

\textbf{Real output:}

\Needspace{82.4pt}
\begin{lstlisting}
D6 for-in: S i
D6 for-in: U y
D6 for-in: Y i
D6 for-of: Siyah 10
D6 for-of: Uyan 7
D6 for-of: Yine 9
D6 assignment targets: {"title":"Siyah","rating":10}
D6 catch: No song selected
\end{lstlisting}

\textbf{What the experiment shows:} The JavaScript for-in loop receives
property-name strings. Its array pattern consumes the first two
characters of each title, producing S/i, U/y and Y/i. The for-of loop
over Object.entries receives actual title/rating pairs. This is a
concrete case where legal JavaScript syntax produces the wrong meaning.
TypeScript rejects the for-in pattern outright, as the negative
experiment below shows. The final two examples assign into existing
object properties and select message in a catch binding. Assignment
patterns can update property targets; declarations introduce bindings. A
catch pattern still assumes the thrown value supports that shape, so
arbitrary thrown null/undefined values need validation instead.

\subsection{D7: Extract tuple/object types and connect them to generic
React
props}\label{d7-extract-tupleobject-types-and-connect-them-to-generic-react-props}

\Needspace{82.4pt}
\begin{lstlisting}[language=TSX]
import type { ReactNode } from 'react'
import { renderToStaticMarkup } from 'react-dom/server'

// These helpers extract types from tuples, not values from runtime arrays.
type First<T extends readonly unknown[]> =
  T extends readonly [infer Head, ...unknown[]] ? Head : never
type Rest<T extends readonly unknown[]> =
  T extends readonly [unknown, ...infer Tail] ? Tail : []
type ValueOf<T> = T extends { value: infer Value } ? Value : never

const queue = ['Siyah', 'Uyan', 'Yine'] as const
const first: First<typeof queue> = queue[0]
const rest: Rest<typeof queue> = ['Uyan', 'Yine']
console.log('D7 First/Rest:', first, JSON.stringify(rest))

type Song = { title: string; rating: number }
function makeSong(title: string, rating: number): Song {
  return { title, rating }
}
const args: Parameters<typeof makeSong> = ['Siyah', 10]
const { title, rating } = makeSong(...args)
const selectedRating: ValueOf<{ value: number }> = rating
console.log('D7 Parameters/ValueOf:', title, selectedRating)

function SongList<T>({ items, children }: {
  items: readonly T[]
  children: (item: T) => ReactNode
}) {
  return <ul>{items.map(item => children(item))}</ul>
}
const songs: Song[] = [
  { title: 'Siyah', rating: 10 },
  { title: 'Uyan', rating: 7 },
  { title: 'Yine', rating: 9 },
]
console.log('D7 generic component:', renderToStaticMarkup(
  <SongList items={songs}>{song => <li key={song.title}>{song.title}: {song.rating}</li>}</SongList>
))
\end{lstlisting}

\Needspace{9\baselineskip}

\textbf{Real output:}

\Needspace{51.5pt}
\begin{lstlisting}
D7 First/Rest: Siyah ["Uyan","Yine"]
D7 Parameters/ValueOf: Siyah 10
D7 generic component: <ul><li>Siyah: 10</li><li>Uyan: 7</li><li>Yine: 9</li></ul>
\end{lstlisting}

First and Rest operate on tuple types at compile time; they do not pull
values out of a runtime array. First is the literal type "Siyah", and
Rest is the tuple {[}"Uyan", "Yine"{]}. This particular Rest helper
infers a mutable tail, even though the input tuple is readonly.
First\textless string{[}{]}\textgreater{} becomes never because an
ordinary array is not guaranteed to have a head. A separate runtime
guard/default is the appropriate way to handle a possibly empty array.

ValueOf uses infer to extract the type of an object\textquotesingle s
value field. Parameters produces the argument tuple {[}title: string,
rating: number{]}, which checks the spread call. These are type-level
structural operations; their annotations disappear at runtime.
\href{https://www.typescriptlang.org/docs/handbook/2/conditional-types.html\#inferring-within-conditional-types}{TypeScript
conditional inference} and
\href{https://www.typescriptlang.org/docs/handbook/utility-types.html\#parameterstype}{Parameters}
describe those type tools.

SongList infers T from its items. The function-valued children prop
therefore receives a Song, and the component calls it once per item to
produce the list. The title/rating fields are checked without an
explicit any, a cast or an additional suppression. This connects generic
typing to the function-children contract from P3, rather than confusing
a type helper with runtime destructuring.

\Needspace{9\baselineskip}

\textbf{Real negative compiler results:}

\Needspace{61.8pt}
\begin{lstlisting}
_bonus_probe.tsx(6,7): error TS2322: Type '"Siyah"' is not assignable to type 'never'.
_bonus_probe.tsx(8,56): error TS2322: Type 'string' is not assignable to type 'number'.
_bonus_probe.tsx(16,22): error TS2339: Property 'duration' does not exist on type '{ title: string; rating: number; }'.
_bonus_probe.tsx(18,12): error TS2491: The left-hand side of a 'for...in' statement cannot be a destructuring pattern.
\end{lstlisting}

The first error verifies that an arbitrary string cannot inhabit the
never result of First\textless string{[}{]}\textgreater. The second
rejects a string rating in Parameters. The third rejects a field absent
from the inferred item type. The fourth is the TypeScript restriction on
the JavaScript for-in experiment. The sources are saved in
bonus-compiler-probe.txt. To reproduce without leaving errors in the
project:

\Needspace{61.8pt}
\begin{lstlisting}[language=sh]
cp bonus-compiler-probe.txt _bonus_probe.tsx
node node_modules/typescript/bin/tsc --noEmit --strict --noUncheckedIndexedAccess --skipLibCheck --target ES2022 --module NodeNext --moduleResolution NodeNext --jsx react-jsx _bonus_probe.tsx --pretty false > bonus-compiler-errors.txt
rm _bonus_probe.tsx
npm run check
\end{lstlisting}

The explicit negative command should fail, while the final ordinary
project check should pass. The original T5 directive remains the
project\textquotesingle s only suppression.

\section{Bonus 3: Real code}\label{bonus-3-real-code}

\subsection{Case 1: The state and dispatch returned by
useReducer}\label{case-1-the-state-and-dispatch-returned-by-usereducer}

\textbf{Source:} Assignment 3 starter (sample Task Manager),
\passthrough{\lstinline!src/App.tsx!}, line 37, inside the
\passthrough{\lstinline!App!} component.

\textbf{Original code:}

\Needspace{30.9pt}
\begin{lstlisting}[language=TSX]
const [tasks, dispatch] = useReducer(taskReducer, [])
\end{lstlisting}

\textbf{Form:} Form 1 - taking values by position from a
function\textquotesingle s returned array.

\Needspace{9\baselineskip}

\textbf{Without destructuring:}

\Needspace{51.5pt}
\begin{lstlisting}[language=TSX]
const reducerResult = useReducer(taskReducer, [])
const tasks = reducerResult[0]
const dispatch = reducerResult[1]
\end{lstlisting}

\Needspace{9\baselineskip}

\textbf{Explanation:}

useReducer returns two values in a fixed order. The first value is the
current state, which is the task list in this starter. The second value
is the dispatch function used to send actions to the reducer.

The array pattern assigns these values to tasks and dispatch by
position. The variable names do not select the values: changing a local
name would not change which position it receives.

Both versions call useReducer exactly once, at the top level of the
component. The version without destructuring stores the returned array
first and then reads its two positions.

\textbf{Readability assessment:}

Destructuring makes this example easier to read because the two related
results are named together. It also avoids an extra variable and
repeated numeric indexes.

\subsection{Case 2: Destructuring the TaskRow
props}\label{case-2-destructuring-the-taskrow-props}

\textbf{Source:} Assignment 3 starter (sample Task Manager),
\passthrough{\lstinline!src/TaskRow.tsx!}, lines 5-14.

\textbf{Original code:}

\Needspace{82.4pt}
\begin{lstlisting}[language=TSX]
export default function TaskRow({ task, onNext, onDelete }: TaskRowProps) {
  return (
    <li style={{ display: 'flex', gap: 8, alignItems: 'center' }}>
      <span style={{ flex: 1 }}>{task.title}</span>
      <em>{task.status}</em>
      <button onClick={onNext}>Next</button>
      <button onClick={onDelete}>Delete</button>
    </li>
  )
}
\end{lstlisting}

\textbf{Form:} Form 5 combined with Form 11 - selecting named properties
inside a function parameter.

\Needspace{9\baselineskip}

\textbf{Without destructuring:}

\Needspace{82.4pt}
\begin{lstlisting}[language=TSX]
export default function TaskRow(props: TaskRowProps) {
  const task = props.task
  const onNext = props.onNext
  const onDelete = props.onDelete

  return (
    <li style={{ display: 'flex', gap: 8, alignItems: 'center' }}>
      <span style={{ flex: 1 }}>{task.title}</span>
      <em>{task.status}</em>
      <button onClick={onNext}>Next</button>
      <button onClick={onDelete}>Delete</button>
    </li>
  )
}
\end{lstlisting}

\Needspace{9\baselineskip}

\textbf{Explanation:}

TaskRow receives one props object. The parameter pattern selects its
task, onNext and onDelete properties. TaskRowProps describes that whole
object; the annotation comes after the pattern.

task is the task object. onNext and onDelete are callback props. Passing
these functions does not call them immediately: the buttons use them as
click handlers.

\textbf{Props and children:}

The starter renders TaskRow as a self-closing component in
\passthrough{\lstinline!src/App.tsx!}, lines 25-30. It passes task,
onNext and onDelete, without nested content. TaskRowProps does not
declare children.

The words Next and Delete are children of the button elements created
inside TaskRow. They are not children passed into TaskRow by its caller.
The key supplied in App is React\textquotesingle s special list
identifier, not an ordinary readable TaskRow prop.

\textbf{Readability assessment:}

The parameter pattern makes the component\textquotesingle s three inputs
visible together. This is useful because the list is short and each
selected name is used directly in the JSX.

\subsection{Case 3: Renaming MessageChannel ports in
esbuild}\label{case-3-renaming-messagechannel-ports-in-esbuild}

\textbf{Source:} Installed dependency \passthrough{\lstinline!esbuild!},
version \passthrough{\lstinline!0.28.2!},
\passthrough{\lstinline!node\_modules/esbuild/lib/main.js!}, line 2394.

\textbf{Original code:}

\Needspace{30.9pt}
\begin{lstlisting}[language=JavaScript]
let { port1: mainPort, port2: workerPort } = new worker_threads2.MessageChannel();
\end{lstlisting}

\textbf{Form:} Form 6 - renaming properties while destructuring.

\Needspace{9\baselineskip}

\textbf{Without destructuring:}

\Needspace{51.5pt}
\begin{lstlisting}[language=JavaScript]
let channel = new worker_threads2.MessageChannel();
let mainPort = channel.port1;
let workerPort = channel.port2;
\end{lstlisting}

\Needspace{9\baselineskip}

\textbf{Explanation:}

The pattern selects the properties named port1 and port2. Their values
are assigned to local variables named mainPort and workerPort.

The names on the left of the colons belong to the source object. The
names on the right are the local bindings. This statement does not
create local variables named port1 or port2.

The version without destructuring creates the channel once and then
reads its two properties. Calling the constructor separately for each
property would create two different channels and would change the
behaviour.

\textbf{Readability assessment:}

The renames are useful here because mainPort and workerPort describe the
roles of the two endpoints more clearly than their numbered property
names.

\subsection{Case 4: A nested worker reply pattern that is harder to
read}\label{case-4-a-nested-worker-reply-pattern-that-is-harder-to-read}

\textbf{Source:} Installed dependency \passthrough{\lstinline!esbuild!},
version \passthrough{\lstinline!0.28.2!},
\passthrough{\lstinline!node\_modules/esbuild/lib/main.js!}, line 2436,
inside runCallSync.

\textbf{Original code:}

\Needspace{30.9pt}
\begin{lstlisting}[language=JavaScript]
let { message: { id: id2, resolve, reject, properties } } = worker_threads2.receiveMessageOnPort(mainPort);
\end{lstlisting}

\textbf{Form:} Form 9 combined with Form 6 - a nested object pattern
containing a rename.

\Needspace{9\baselineskip}

\textbf{Without destructuring:}

\Needspace{82.4pt}
\begin{lstlisting}[language=JavaScript]
let packet = worker_threads2.receiveMessageOnPort(mainPort);
let message = packet.message;
let id2 = message.id;
let resolve = message.resolve;
let reject = message.reject;
let properties = message.properties;
\end{lstlisting}

\Needspace{9\baselineskip}

\textbf{Explanation:}

The returned packet contains a message property. The nested pattern
reads id, resolve, reject and properties from that message object. The
id property is renamed to id2.

In the surrounding source, runCallSync already has a local variable
named id for the outgoing request. The next line compares that id with
id2 from the reply.

The original pattern does not create a local variable named message. It
uses message only as a step into the nested object. It also supplies no
default: an undefined packet or a missing message object would cause the
extraction to fail.

\textbf{Readability assessment:}

This is the case where the compact destructuring makes the code harder
to read. One long statement combines a function call, two object levels,
four selected fields and a rename.

Separating the packet and message steps exposes the reply structure and
provides useful places to inspect values while debugging. The longer
version also makes it easier to see which object supplies each field. It
preserves the single receiveMessageOnPort call and does not add fallback
behaviour.

\subsection{Case 5: Destructuring array entries in
Prettier}\label{case-5-destructuring-array-entries-in-prettier}

\textbf{Source:} Prettier version \passthrough{\lstinline!3.5.3!},
\passthrough{\lstinline!src/language-js/print/jsx.js!}, lines 382-388.

\href{https://github.com/prettier/prettier/blob/3.5.3/src/language-js/print/jsx.js\#L382-L388}{View
the original source}.

\textbf{Original code:}

\Needspace{82.4pt}
\begin{lstlisting}[language=JavaScript]
for (const [i, word] of words.entries()) {
  if (i % 2 === 1) {
    pushLine(line);
  } else {
    push(word);
  }
}
\end{lstlisting}

\textbf{Form:} Form 12 - destructuring pairs inside a for...of loop.
This example uses Array.prototype.entries rather than Object.entries or
Map.

\Needspace{9\baselineskip}

\textbf{Without destructuring:}

\Needspace{82.4pt}
\begin{lstlisting}[language=JavaScript]
for (const entry of words.entries()) {
  const i = entry[0];
  const word = entry[1];

  if (i % 2 === 1) {
    pushLine(line);
  } else {
    push(word);
  }
}
\end{lstlisting}

\Needspace{9\baselineskip}

\textbf{Explanation:}

words.entries() supplies an iterator of index/value pairs. The array
pattern assigns the index to i and the corresponding value to word on
each iteration.

The loop uses the index\textquotesingle s parity to choose between
pushLine(line) and push(word). The version without destructuring reads
the same two positions from each entry and leaves the loop body
unchanged.

These are array indexes, not property names obtained from for...in. This
code processes JSX formatting inside Prettier; it is not a React
component receiving a children prop.

\textbf{Readability assessment:}

Destructuring works well here because both parts of every pair are
needed. Naming them in the loop header makes their relationship clearer
than repeatedly reading entry{[}0{]} and entry{[}1{]}.

\subsection{Source and version notes}\label{source-and-version-notes}

Cases 1 and 2 come from the instructor\textquotesingle s original
Assignment 3 starter, rather than treating my completed assignment as
somebody else\textquotesingle s project.

Cases 3 and 4 were checked in the installed esbuild dependency of the
clean-clone verification copy. This project\textquotesingle s
package-lock.json records esbuild version 0.28.2. The line numbers refer
to that installed version.

Case 5 links to a specific Prettier release so the file and line
references can be checked against the version analysed here.

These five cases cover the three requested source groups: the Task
Manager starter, an installed library and a separate open-source
project.

\section{Bonus 4: Git repository}\label{bonus-4-git-repository}

https://github.com/mertsengel125-coder/assignment4-destructuring-field-guide

The repository is public and provides read access to the source.
node\_modules is excluded by .gitignore. Reproduce the project with:

\Needspace{72.1pt}
\begin{lstlisting}[language=sh]
git clone https://github.com/mertsengel125-coder/assignment4-destructuring-field-guide.git
cd assignment4-destructuring-field-guide
npm install
npm run all
npm run check
\end{lstlisting}

The repository history contains at least six commits, with work recorded
on 19 and 21 November 2026. The commits cover the required guide,
runnable examples, depth investigations, LaTeX/PDF, clean-clone
verification and real-code analyses. The history spans two calendar
dates. All commit dates are real. The instructor can inspect the actual
history to assess the requirement for work over several days.

\section{Submission checks}\label{submission-checks}

\begin{itemize}
\tightlist
\item
  Complete and understand the examples and explanations, including the
  difference between rest and spread.
\item
  Regenerate output.txt from the final npm run all and check each README
  output block.
\item
  Review all five personal log entries and keep each entry consistent
  with the actual experiment and output.
\item
  Keep the correct name/student number at the top of each source file.
\item
  Keep strict settings and only the single T5 suppression.
\item
  Zip one assignment4 folder without node\_modules or .git. Upload to
  Moodle by its announced deadline.
\item
  Bonuses 1, 2, 3 and 4 are claimed above. Check the actual repository
  history and ensure the final code, recorded outputs, Markdown guide
  and LaTeX/PDF agree.
\end{itemize}

\section{My own log}\label{my-own-log}

\subsection{1. Comparing my prediction with the playlist
summary}\label{comparing-my-prediction-with-the-playlist-summary}

Before running the code, I expected the total duration to be 740 seconds
and the average rating to be 8.

\textbf{Command:}

\Needspace{30.9pt}
\begin{lstlisting}[language=sh]
node js/arrays.mjs
\end{lstlisting}

\textbf{Actual output:}

\Needspace{30.9pt}
\begin{lstlisting}
F1: 777 8.666666666666666
\end{lstlisting}

\Needspace{9\baselineskip}

\textbf{Explanation:}

The total duration is 285 + 231 + 261 = 777 seconds. The average rating
is (10 + 7 + 9) / 3 = 26 / 3, which is approximately 8.67.

The function returns an array containing the total duration first and
the average rating second. Array destructuring assigns these values to
playlistSeconds and averageRating in that order. The positions determine
the assignments, not the variable names.

\subsection{2. Defaults for undefined and
null}\label{defaults-for-undefined-and-null}

\textbf{What I actually tried:} I ran
\passthrough{\lstinline!node js/pitfalls.mjs!} and compared the missing,
undefined, null and zero cases.

\textbf{Actual output:}

\Needspace{61.8pt}
\begin{lstlisting}
Pitfall 2 missing: 35
Pitfall 2 undefined: 35
Pitfall 2 null: null
Pitfall 2 zero: 0
\end{lstlisting}

\textbf{Observed result:} The missing and undefined cases printed 35.
The null and zero cases kept their supplied values.

\Needspace{9\baselineskip}

\textbf{Explanation:}

In \passthrough{\lstinline!const \{ volume = 35 \} = settings!}, the
default is used only when the extracted value is
\passthrough{\lstinline!undefined!}. A missing property also produces
\passthrough{\lstinline!undefined!}, so both of those cases use 35. An
explicit \passthrough{\lstinline!null!} remains null, and 0 remains 0.
Zero is a valid volume value, so a destructuring default does not
replace it. If both null and undefined should use a fallback, that
requires a separate rule such as \passthrough{\lstinline!value ?? 35!}.

\subsection{3. Assignment parentheses}\label{assignment-parentheses}

\textbf{What I actually tried:} I ran the invalid and parenthesized
assignment experiments in
\passthrough{\lstinline!node js/pitfalls.mjs!}.

\textbf{Actual output:}

\Needspace{41.2pt}
\begin{lstlisting}
Pitfall 3 invalid: SyntaxError: Unexpected token '='
Pitfall 3 fixed: Siyah
\end{lstlisting}

\textbf{Observed result:} The invalid attempt produced a SyntaxError.
The parenthesized assignment printed Siyah. The invalid attempt was
isolated and caught, so the rest of the file continued running.

\Needspace{9\baselineskip}

\textbf{Explanation:}

At the beginning of a statement, JavaScript reads
\passthrough{\lstinline!\{ title \}!} as a block rather than an object
assignment pattern. The following \passthrough{\lstinline!=!} therefore
causes a SyntaxError. Writing
\passthrough{\lstinline!(\{ title \} = \{ title: 'Siyah' \})!} puts the
assignment inside an expression, so the pattern assigns the title
property to the existing title variable. The invalid attempt is parsed
separately with \passthrough{\lstinline!new Function!} and caught;
otherwise invalid syntax in the main file would prevent that file from
running.

\subsection{4. A shared nested object}\label{a-shared-nested-object}

\textbf{What I actually tried:} I ran the rest-copy experiment in
\passthrough{\lstinline!node js/pitfalls.mjs!}, which assigns 7 to
copy.stats.rating.

\textbf{Actual output:}

\Needspace{51.5pt}
\begin{lstlisting}
Pitfall 4 original rating: 7
Pitfall 4 different outer object: true
Pitfall 4 shared stats: true
\end{lstlisting}

\textbf{Observed result:} The original rating became 7. The outer
objects were different, while both stats properties referred to the same
nested object.

\Needspace{9\baselineskip}

\textbf{Explanation:}

Object rest creates a new outer object, but it copies the reference to
the nested stats object. As a result,
\passthrough{\lstinline"copy !== song"} is true, while
\passthrough{\lstinline!copy.stats === song.stats!} is also true. Both
stats properties refer to the same nested object. Assigning 7 to
\passthrough{\lstinline!copy.stats.rating!} therefore makes
\passthrough{\lstinline!song.stats.rating!} show 7 too. This is a
shallow copy, not a deep copy of the nested data.

\subsection{5. A function passed as
children}\label{a-function-passed-as-children}

\textbf{What I actually tried:} I ran
\passthrough{\lstinline!npx tsx js/props.jsx!} and compared the raw
function child with the version where PlaylistSummary calls the
function.

\textbf{Actual output:}

\Needspace{51.5pt}
\begin{lstlisting}
P3 raw function: <div></div>
P3 diagnostic: Functions are not valid as a React child. This may happen if you return renderTrackCount instead of <renderTrackCount /> from render. Or maybe you meant to call this function rather than return it.
P3 called function: <div><strong>3 tracks</strong></div>
\end{lstlisting}

\textbf{Observed result:} Passing the function itself to Frame produced
an empty div and a diagnostic. When PlaylistSummary called the function
with the count 3, the returned strong element appeared in the HTML.

\Needspace{9\baselineskip}

\textbf{Explanation:}

Passing a function as a prop does not call it. Frame places its children
directly inside a div, and an ordinary function is not a renderable
React child. PlaylistSummary uses a different contract: its children
prop is a function, and the component explicitly calls children(count).
The returned React element can then be rendered. In TypeScript,
children?: ReactNode rejects an ordinary function, while children:
(count: number) =\textgreater{} ReactNode describes the function
contract used by PlaylistSummary.

\end{document}
